\documentclass{article}

\begin{document}
	
	SUPPORT VECTOR MACHINES — ADVANCED EXAM WITH ANSWERS
	
	\bigskip
	==================================================
	
	PARAGRAPH 1 — INTRODUCTION TO SVM
	
	\begin{enumerate}
		\item Which learning paradigm best describes SVM?
		\begin{enumerate}
			\item Unsupervised learning
			\item Reinforcement learning
			\item Supervised learning
			\item Self-supervised learning
		\end{enumerate}
		
		\item The original formulation of SVM targets:
		\begin{enumerate}
			\item Density estimation
			\item Binary classification
			\item Multiclass classification
			\item Regression only
		\end{enumerate}
		
		\item SVM theory is primarily based on:
		\begin{enumerate}
			\item Empirical risk minimization
			\item Structural risk minimization
			\item Maximum likelihood estimation
			\item Bayesian inference
		\end{enumerate}
		
		\item A key reason SVMs generalize well is:
		\begin{enumerate}
			\item Large training sets
			\item Margin maximization
			\item Feature normalization
			\item Kernel complexity
		\end{enumerate}
		
		\item Which property distinguishes SVMs from k-NN?
		\begin{enumerate}
			\item Non-parametric nature
			\item Margin-based optimization
			\item Distance-based decisions
			\item Instance-based learning
		\end{enumerate}
		
		\item In their basic form, SVMs assume data are:
		\begin{enumerate}
			\item Noisy
			\item Perfectly separable
			\item Sequential
			\item Probabilistic
		\end{enumerate}
		
		\item Which quantity is explicitly optimized during training?
		\begin{enumerate}
			\item Classification accuracy
			\item Margin width
			\item Likelihood
			\item Entropy
		\end{enumerate}
		
		\item Which statement is true for classical SVMs?
		\begin{enumerate}
			\item They minimize squared error
			\item They produce probabilistic outputs
			\item They rely on convex optimization
			\item They require gradient descent
		\end{enumerate}
		
		\item Compared to neural networks, SVMs:
		\begin{enumerate}
			\item Always require more data
			\item Have a unique global optimum
			\item Are always faster to train
			\item Learn hierarchical features
		\end{enumerate}
		
		\item SVMs are particularly effective when:
		\begin{enumerate}
			\item Number of features is very small
			\item Data dimensionality is high
			\item Classes overlap heavily
			\item Labels are uncertain
		\end{enumerate}
	\end{enumerate}
	
	Correct answers: 1C, 2B, 3B, 4B, 5B, 6B, 7B, 8C, 9B, 10B
	
	\bigskip
	==================================================
	
	PARAGRAPH 2 — TRAINING DATA
	
	\begin{enumerate}
		\item Each SVM training sample is represented as:
		\begin{enumerate}
			\item Feature vector only
			\item Label only
			\item Feature vector and label
			\item Probability and label
		\end{enumerate}
		
		\item The label space in classical SVMs is:
		\begin{enumerate}
			\item Continuous
			\item Multiclass
			\item Binary
			\item Ordinal
		\end{enumerate}
		
		\item Why are labels encoded as minus one and plus one?
		\begin{enumerate}
			\item Numerical convenience
			\item Probabilistic interpretation
			\item Kernel compatibility
			\item Data normalization
		\end{enumerate}
		
		\item Feature dimensionality affects:
		\begin{enumerate}
			\item Margin definition
			\item Kernel computation
			\item Decision boundary complexity
			\item All of the above
		\end{enumerate}
		
		\item Training data must be:
		\begin{enumerate}
			\item Ordered
			\item Fully labeled
			\item Balanced
			\item Normalized
		\end{enumerate}
		
		\item If labels are noisy, SVM performance:
		\begin{enumerate}
			\item Always improves
			\item Becomes undefined
			\item Requires soft margin
			\item Requires kernel removal
		\end{enumerate}
		
		\item The index i in training pairs denotes:
		\begin{enumerate}
			\item Feature index
			\item Dimension index
			\item Sample index
			\item Class index
		\end{enumerate}
		
		\item What happens if training labels are incorrect?
		\begin{enumerate}
			\item Margin increases
			\item Optimization fails
			\item Slack variables become active
			\item Kernel degenerates
		\end{enumerate}
		
		\item SVMs can handle missing labels:
		\begin{enumerate}
			\item Always
			\item With modification
			\item Never
			\item Only with kernels
		\end{enumerate}
		
		\item Feature scaling is important because:
		\begin{enumerate}
			\item It changes labels
			\item It affects margin geometry
			\item It removes support vectors
			\item It avoids overfitting
		\end{enumerate}
	\end{enumerate}
	
	Correct answers: 1C, 2C, 3A, 4D, 5B, 6C, 7C, 8C, 9B, 10B
	
	\bigskip
	==================================================
	
	PARAGRAPH 3 — HYPERPLANE AND DECISION FUNCTION
	
	\begin{enumerate}
		\item The separating hyperplane is defined as:
		\begin{enumerate}
			\item w + x + b = 0
			\item w transpose x plus b equals zero
			\item x transpose x equals one
			\item w equals zero
		\end{enumerate}
		
		\item The vector w is:
		\begin{enumerate}
			\item Parallel to the hyperplane
			\item Orthogonal to the hyperplane
			\item Tangent to the margin
			\item Independent of data
		\end{enumerate}
		
		\item The bias term b determines:
		\begin{enumerate}
			\item Orientation
			\item Offset
			\item Margin width
			\item Regularization
		\end{enumerate}
		
		\item Decision boundary corresponds to:
		\begin{enumerate}
			\item Maximum margin
			\item Zero decision function
			\item Maximum likelihood
			\item Kernel output
		\end{enumerate}
		
		\item Classification depends on:
		\begin{enumerate}
			\item Absolute value of w
			\item Sign of decision function
			\item Kernel type
			\item Number of samples
		\end{enumerate}
		
		\item Points exactly on the boundary satisfy:
		\begin{enumerate}
			\item Positive score
			\item Negative score
			\item Zero score
			\item Unit score
		\end{enumerate}
		
		\item Changing b while fixing w will:
		\begin{enumerate}
			\item Rotate the hyperplane
			\item Shift the hyperplane
			\item Change margin size
			\item Change kernel
		\end{enumerate}
		
		\item Scaling w and b by the same factor:
		\begin{enumerate}
			\item Changes the boundary
			\item Preserves the boundary
			\item Eliminates margin
			\item Changes labels
		\end{enumerate}
		
		\item Decision function values can be interpreted as:
		\begin{enumerate}
			\item Probabilities
			\item Distances up to scale
			\item Likelihoods
			\item Errors
		\end{enumerate}
		
		\item A larger absolute decision value implies:
		\begin{enumerate}
			\item Lower confidence
			\item Higher confidence
			\item Misclassification
			\item Kernel failure
		\end{enumerate}
	\end{enumerate}
	
	Correct answers: 1B, 2B, 3B, 4B, 5B, 6C, 7B, 8B, 9B, 10B
	
	\bigskip
	==================================================
	
	PARAGRAPH 4 — MARGIN AND SUPPORT VECTORS
	
	\begin{enumerate}
		\item The margin is defined as:
		\begin{enumerate}
			\item Distance to origin
			\item Distance between centroids
			\item Distance between margin hyperplanes
			\item Distance between all samples
		\end{enumerate}
		
		\item Maximizing margin corresponds to:
		\begin{enumerate}
			\item Minimizing w norm
			\item Maximizing likelihood
			\item Minimizing error
			\item Maximizing entropy
		\end{enumerate}
		
		\item Support vectors are points that:
		\begin{enumerate}
			\item Are misclassified
			\item Lie on the margin
			\item Are far from boundary
			\item Are randomly chosen
		\end{enumerate}
		
		\item Removing non-support vectors:
		\begin{enumerate}
			\item Changes solution
			\item Has no effect
			\item Increases margin
			\item Breaks optimization
		\end{enumerate}
		
		\item Margin width depends on:
		\begin{enumerate}
			\item All samples
			\item Support vectors only
			\item Kernel type only
			\item Dimensionality only
		\end{enumerate}
		
		\item A small margin usually implies:
		\begin{enumerate}
			\item Better generalization
			\item Overfitting
			\item Underfitting
			\item No solution
		\end{enumerate}
		
		\item Support vectors have:
		\begin{enumerate}
			\item Zero multipliers
			\item Positive multipliers
			\item Negative multipliers
			\item No multipliers
		\end{enumerate}
		
		\item Margin maximization is a form of:
		\begin{enumerate}
			\item Empirical risk minimization
			\item Structural risk minimization
			\item Heuristic optimization
			\item Bayesian inference
		\end{enumerate}
		
		\item Increasing dimensionality generally:
		\begin{enumerate}
			\item Decreases margin
			\item Has no effect
			\item Increases separability
			\item Eliminates support vectors
		\end{enumerate}
		
		\item The margin is invariant to:
		\begin{enumerate}
			\item Scaling of w and b
			\item Feature normalization
			\item Kernel choice
			\item Slack variables
		\end{enumerate}
	\end{enumerate}
	
	Correct answers: 1C, 2A, 3B, 4B, 5B, 6B, 7B, 8B, 9C, 10A
	
	\bigskip
	==================================================
	
	PARAGRAPH 5 — SOFT MARGIN AND PARAMETER C
	
	\begin{enumerate}
		\item Slack variables represent:
		\begin{enumerate}
			\item Feature noise
			\item Margin violations
			\item Kernel error
			\item Scaling error
		\end{enumerate}
		
		\item The parameter C controls:
		\begin{enumerate}
			\item Kernel smoothness
			\item Margin-error tradeoff
			\item Feature scaling
			\item Dimensionality
		\end{enumerate}
		
		\item A large C leads to:
		\begin{enumerate}
			\item Wider margins
			\item More regularization
			\item Fewer training errors
			\item Smaller model
		\end{enumerate}
		
		\item A small C encourages:
		\begin{enumerate}
			\item Hard-margin behavior
			\item Wider margin
			\item Perfect classification
			\item Kernel removal
		\end{enumerate}
		
		\item Soft-margin SVM is useful when:
		\begin{enumerate}
			\item Data are noise-free
			\item Classes overlap
			\item Features are binary
			\item Kernels are linear
		\end{enumerate}
		
		\item If C tends to infinity:
		\begin{enumerate}
			\item Margin dominates
			\item Errors dominate
			\item Hard-margin is recovered
			\item Optimization fails
		\end{enumerate}
		
		\item If C tends to zero:
		\begin{enumerate}
			\item Margin ignored
			\item Errors ignored
			\item Margin dominates
			\item No solution exists
		\end{enumerate}
		
		\item Slack variables are constrained to be:
		\begin{enumerate}
			\item Negative
			\item Zero
			\item Non-negative
			\item Binary
		\end{enumerate}
		
		\item Soft-margin optimization remains:
		\begin{enumerate}
			\item Non-convex
			\item Convex
			\item Discrete
			\item Probabilistic
		\end{enumerate}
		
		\item Increasing C generally:
		\begin{enumerate}
			\item Reduces overfitting
			\item Increases overfitting
			\item Has no effect
			\item Removes kernels
		\end{enumerate}
	\end{enumerate}
	
	Correct answers: 1B, 2B, 3C, 4B, 5B, 6C, 7C, 8C, 9B, 10B
	
	\bigskip
	==================================================
	
	PARAGRAPH 6 — NONLINEAR SVM AND KERNELS
	
	\begin{enumerate}
		\item Kernels allow SVMs to:
		\begin{enumerate}
			\item Reduce dimensions
			\item Operate in implicit feature spaces
			\item Avoid optimization
			\item Remove slack variables
		\end{enumerate}
		
		\item The kernel trick replaces:
		\begin{enumerate}
			\item Optimization
			\item Dot products
			\item Regularization
			\item Labels
		\end{enumerate}
		
		\item Which kernel maps to infinite dimensions?
		\begin{enumerate}
			\item Linear
			\item Polynomial
			\item Gaussian
			\item Sigmoid
		\end{enumerate}
		
		\item Polynomial kernels control complexity via:
		\begin{enumerate}
			\item Degree
			\item Bias only
			\item Dimensionality
			\item Labels
		\end{enumerate}
		
		\item Kernel choice affects:
		\begin{enumerate}
			\item Optimization convexity
			\item Feature space geometry
			\item Label encoding
			\item Slack constraints
		\end{enumerate}
		
		\item A linear kernel implies:
		\begin{enumerate}
			\item Nonlinear boundary
			\item Original feature space
			\item Infinite dimensions
			\item Probabilistic output
		\end{enumerate}
		
		\item Overly complex kernels may cause:
		\begin{enumerate}
			\item Underfitting
			\item Overfitting
			\item No solution
			\item Kernel failure
		\end{enumerate}
		
		\item Kernel parameters must be:
		\begin{enumerate}
			\item Fixed arbitrarily
			\item Learned analytically
			\item Tuned via validation
			\item Ignored
		\end{enumerate}
		
		\item The kernel matrix must be:
		\begin{enumerate}
			\item Symmetric positive semi-definite
			\item Diagonal
			\item Invertible
			\item Sparse
		\end{enumerate}
		
		\item Nonlinear SVMs remain:
		\begin{enumerate}
			\item Linear classifiers
			\item Convex optimization problems
			\item Probabilistic models
			\item Neural networks
		\end{enumerate}
	\end{enumerate}
	
	Correct answers: 1B, 2B, 3C, 4A, 5B, 6B, 7B, 8C, 9A, 10B
	
	\bigskip
	==================================================
	
	PARAGRAPH 7 — MULTICLASS SVM
	
	\begin{enumerate}
		\item One-vs-all strategy trains:
		\begin{enumerate}
			\item One classifier
			\item K classifiers
			\item K(K-1) classifiers
			\item Two classifiers
		\end{enumerate}
		
		\item Each classifier separates:
		\begin{enumerate}
			\item Two classes
			\item One class vs all others
			\item All classes jointly
			\item Random subsets
		\end{enumerate}
		
		\item Final prediction is based on:
		\begin{enumerate}
			\item Minimum margin
			\item Maximum decision score
			\item Random choice
			\item Kernel output
		\end{enumerate}
		
		\item One-vs-one strategy requires:
		\begin{enumerate}
			\item K classifiers
			\item K-1 classifiers
			\item K(K-1)/2 classifiers
			\item 2K classifiers
		\end{enumerate}
		
		\item Multiclass SVMs are:
		\begin{enumerate}
			\item Exact extensions
			\item Heuristic decompositions
			\item Unsupervised
			\item Kernel-free
		\end{enumerate}
		
		\item Which strategy scales better with classes?
		\begin{enumerate}
			\item One-vs-all
			\item One-vs-one
			\item Both equally
			\item Neither
		\end{enumerate}
		
		\item In practice, multiclass SVM performance depends on:
		\begin{enumerate}
			\item Kernel choice
			\item Decomposition strategy
			\item Data distribution
			\item All of the above
		\end{enumerate}
		
		\item Ambiguities in multiclass SVM arise when:
		\begin{enumerate}
			\item All scores are negative
			\item Multiple classifiers give high scores
			\item Kernel fails
			\item Margin is zero
		\end{enumerate}
		
		\item One-vs-all may suffer from:
		\begin{enumerate}
			\item Class imbalance
			\item Convexity loss
			\item Kernel invalidity
			\item No optimization
		\end{enumerate}
		
		\item Multiclass SVMs remain:
		\begin{enumerate}
			\item Non-convex
			\item Convex per subproblem
			\item Probabilistic
			\item Unsupervised
		\end{enumerate}
	\end{enumerate}
	
	Correct answers: 1B, 2B, 3B, 4C, 5B, 6A, 7D, 8B, 9A, 10B
	
\end{document}
